\documentclass[10pt]{article}
\usepackage[a4paper,margin=0.66in]{geometry}
\usepackage[T1]{fontenc}
\usepackage[utf8]{inputenc}
\usepackage{booktabs}
\usepackage{array}
\usepackage{hyperref}
\usepackage{inconsolata}
\setlength{\parindent}{0pt}
\setlength{\parskip}{3pt}
\setlength{\tabcolsep}{4pt}

\title{\vspace{-1.35cm}DiCo-NLI: Trial Data Feasibility and Translation Validation}
\author{SemEval-2027 Task Proposal Supplementary Material, Item 2}
\date{\vspace{-0.7cm}}

\begin{document}
\maketitle
\vspace{-0.8cm}

\section{Reviewer Concern and Scope of This Item}

This item responds to the reviewer request for concrete evidence that the
DiCo-NLI trial data can be constructed reliably before the full shared-task
release.  The main multilingual feasibility risk is target-side translation
collapse: some English pairs whose two phrases are distinct may translate to
exactly the same phrase in Spanish or Basque.  Such cases are unusable for
DiCo-NLI if the target-language pair collapses, because the translated item no
longer presents two distinguishable phrase arguments for directional inference.

This item documents the current trial-data construction.  The evidence shows
that (i) the full English source inventory has been grouped and audited, (ii)
the collapse risk is measurable before release, (iii) multiple high-quality
translation models provide recoverable alternatives for many risky cases, and
(iv) the released trial bundle is generated from accepted non-collapsed Spanish
and Basque candidates.

\section{Full English Source Inventory}

We rebuilt the grouped English database from the class-segregated PhrasIS
layout.  The grouping unit is the source phrase pair; all directions, languages,
and participant-track instances derived from the same source pair remain tied to
that unit for splitting and leakage control.  Table~\ref{tab:full-source}
summarizes the resulting English inventory.

\begin{table}[h]
\centering
\small
\begin{tabular}{lrr}
\toprule
Source class & DiCo-NLI label & Pairs \\
\midrule
\texttt{BACK} & \texttt{BACKWARD\_ENTAILMENT} & 636 \\
\texttt{EQUI} & \texttt{EQUIVALENCE} & 686 \\
\texttt{FORW} & \texttt{FORWARD\_ENTAILMENT} & 624 \\
\texttt{OPPO} & \texttt{NEGATIVE\_OTHER} & 19 \\
\texttt{REL} & \texttt{NEGATIVE\_OTHER} & 411 \\
\texttt{SIMI} & \texttt{NEGATIVE\_OTHER} & 945 \\
\texttt{UNR} & \texttt{NEGATIVE\_OTHER} & 6893 \\
\midrule
\textbf{Total} & -- & \textbf{10214} \\
\bottomrule
\end{tabular}
\caption{Full grouped English source database derived from PhrasIS.}
\label{tab:full-source}
\end{table}

The reversible positive component contains 1946 pairs: 686
\texttt{EQUIVALENCE}, 624 forward-entailment, and 636 backward-entailment
examples.  The remaining 8268 pairs are available as \texttt{NEGATIVE\_OTHER};
they will be capped and stratified so that the final task is not dominated by
negative examples.

\section{Translation Collapse Definition}

For each generated target-language candidate we compute a simple but strict
automatic diagnostic:

\begin{quote}
\textbf{collapse}: after Unicode normalization, case-folding, and whitespace
normalization, translated \texttt{text1} is exactly equal to translated
\texttt{text2}.
\end{quote}

Collapse is not the only validation criterion, but it is a necessary rejection
criterion for target-language DiCo-NLI items.  If \emph{South Korean President}
and \emph{S. Korean president} both become the identical Spanish phrase
\emph{presidente surcoreano}, the target-language pair no longer tests the same
fine-grained distinction as the English source.  This is especially frequent for
\texttt{EQUIVALENCE}, where English often expresses the same meaning with
article, abbreviation, morphology, or word-order variation that may naturally
disappear under translation.

The translation prompt used in this experiment requested faithful phrase-pair
translation and explicitly asked the model to preserve lexical or syntactic
distinctions when a natural, relation-preserving alternative exists.  The prompt
also allows identical translations when identity is genuinely the best
translation.  Thus the automatic collapse flag is not treated as model failure:
it is a data-quality signal telling us whether the resulting translated pair is
suitable for the benchmark.

\section{Stratified 500-Pair Feasibility Experiment}

We sampled 500 source pairs from the full grouped English database with a fixed
seed and equal DICO label strata: 125 pairs per label.  The sample contains
enough \texttt{EQUIVALENCE} examples to quantify the translation collapse risk
directly, while also covering the directional and negative labels used in the
task.

\begin{table}[h]
\centering
\small
\begin{tabular}{lr}
\toprule
DiCo-NLI label & Sampled pairs \\
\midrule
\texttt{BACKWARD\_ENTAILMENT} & 125 \\
\texttt{EQUIVALENCE} & 125 \\
\texttt{FORWARD\_ENTAILMENT} & 125 \\
\texttt{NEGATIVE\_OTHER} & 125 \\
\midrule
\textbf{Total} & \textbf{500} \\
\bottomrule
\end{tabular}
\caption{Stratified sample used for the translation-validation experiment.}
\label{tab:sample}
\end{table}

Every sampled pair was translated independently into Spanish and Basque with
two OpenAI models: \texttt{gpt-5.4} and \texttt{gpt-5.6-sol}.  This produced
1000 translation calls per model: 500 Spanish and 500 Basque.  The
\texttt{gpt-5.6-sol} run was executed in 20 persisted chunks of 25 pairs after a
provider-side HTTP 500 interrupted an initial unchunked attempt.  This directly
informs our production plan: all full-dataset translation runs will use
checkpointed chunks rather than a single all-or-nothing process.

\section{Single-Model Collapse Risk}

Table~\ref{tab:collapse-single} reports the exact collapse counts for each
model.  The results confirm the central risk: a substantial fraction of
\texttt{EQUIVALENCE} pairs collapse in at least one target language if only one
translation model is used.

\begin{table}[h]
\centering
\small
\begin{tabular}{lrrrr}
\toprule
Model & ES collapse & EU collapse & Any language & Both languages \\
\midrule
\texttt{gpt-5.4} & 41 & 52 & 67 & 26 \\
\texttt{gpt-5.6-sol} & 28 & 44 & 50 & 22 \\
\bottomrule
\end{tabular}
\caption{Target-side exact-collapse counts over the same 500 source pairs.}
\label{tab:collapse-single}
\end{table}

For \texttt{gpt-5.4}, 65 of the 67 pairs that collapsed in Spanish or Basque
were \texttt{EQUIVALENCE}.  This corresponds to 65/125 = 52.0\% of the sampled
\texttt{EQUIVALENCE} stratum.  If applied naively to the full inventory of 686
\texttt{EQUIVALENCE} pairs, a single-model deletion policy would remove roughly
357 equivalence pairs, leaving about 329.  This is not acceptable as a final
production strategy, but it is useful evidence that the risk is now measurable
and can be controlled before release.

\section{Benefit of Multiple Translation Models}

The two models do not collapse exactly the same pairs.  Table~\ref{tab:overlap}
shows the cross-model overlap by target language.  For Spanish, 19 pairs
collapsed under \texttt{gpt-5.4} but not under \texttt{gpt-5.6-sol}; for Basque,
20 such pairs were recovered.  Conversely, a smaller number collapsed only under
\texttt{gpt-5.6-sol}.  This supports a multi-candidate workflow: generate more
than one translation candidate, reject candidates that collapse, and send the
remaining candidates to bilingual validation.

\begin{table}[h]
\centering
\small
\begin{tabular}{lrrrr}
\toprule
Language & Both collapse & 5.4 only & 5.6 only & Neither \\
\midrule
Spanish & 22 & 19 & 6 & 453 \\
Basque & 32 & 20 & 12 & 436 \\
\bottomrule
\end{tabular}
\caption{Cross-model collapse overlap.  ``5.4 only'' means the pair collapsed
with \texttt{gpt-5.4} but had a non-collapsed \texttt{gpt-5.6-sol} candidate.}
\label{tab:overlap}
\end{table}

The most relevant estimate is the best-available-candidate scenario: for each
language, keep the pair if at least one model produced a non-collapsed candidate
that can then be manually checked.  Under this policy, only 39 of the 500
sampled pairs had unavoidable collapse in Spanish or Basque across both models.
Within the \texttt{EQUIVALENCE} stratum, 86 of 125 pairs had a non-collapsed
option for both target languages, while 39 remained unavoidable exact-collapse
cases.  Projected to the full 686-pair \texttt{EQUIVALENCE} inventory, this
would leave roughly 472 usable equivalence pairs before any additional
translation engines or manual repair.

\begin{table}[h]
\centering
\small
\begin{tabular}{lr}
\toprule
Best-available-candidate metric & Pairs \\
\midrule
Sample pairs with unavoidable collapse in ES or EU & 39 / 500 \\
Sample pairs with a non-collapsed option for both ES and EU & 461 / 500 \\
\texttt{EQUIVALENCE} pairs with unavoidable collapse & 39 / 125 \\
\texttt{EQUIVALENCE} pairs with usable non-collapsed options & 86 / 125 \\
\bottomrule
\end{tabular}
\caption{Estimated effect of selecting the best non-collapsed candidate per
language from the two tested models.}
\label{tab:best}
\end{table}

This result changes the production plan.  We will not translate once and delete
all flagged pairs.  Instead, we will maintain multiple translation candidates
per language, use exact-collapse and source-copy diagnostics for automatic
triage, and preserve all candidate provenance so that bilingual reviewers can
select, correct, or reject items with full information.

The collapse signal is also useful as a quality filter rather than only a loss
mechanism.  If independent strong translation systems reduce an English
equivalence pair to the same target-language unit, the source pair is often a
near-trivial equivalence: the two phrases may differ only by punctuation,
article choice, abbreviation, or a very small edit.  Removing persistent
cross-model collapse cases therefore helps the final benchmark avoid easy
surface-equivalence items and keeps the retained \texttt{EQUIVALENCE} examples
more challenging for systems.

\section{Cost and Operational Feasibility}

The experiment also gives realistic cost and runtime evidence for full-data
production.  Table~\ref{tab:cost} reports the recorded API usage for the
successful 500-pair runs.  The \texttt{gpt-5.6-sol} model was more expensive and
slower, partly because it used reasoning tokens, but it recovered a useful
subset of \texttt{gpt-5.4} collapses.  This makes it appropriate as a second
candidate generator or targeted repair model rather than necessarily as the only
default generator.

\begin{table}[h]
\centering
\small
\begin{tabular}{lrrrr}
\toprule
Model & Input tok. & Output tok. & Reasoning tok. & Est. cost \\
\midrule
\texttt{gpt-5.4} & 198904 & 25156 & 0 & \$0.8746 \\
\texttt{gpt-5.6-sol} & 198904 & 72558 & 44561 & \$3.1713 \\
\bottomrule
\end{tabular}
\caption{Recorded usage for 500 pairs translated into Spanish and Basque.}
\label{tab:cost}
\end{table}

The first unchunked \texttt{gpt-5.6-sol} attempt failed with an HTTP 500 before
the final JSON could be written.  The successful rerun used 20 saved chunks.
This is an important engineering conclusion: final production will include
progress logging, checkpointed chunk outputs, and resumable merging before any
large translation job is launched.

\section{Validation Protocol for Trial and Full Release}

The final trial and evaluation data will be produced by the following protocol.

\begin{enumerate}
\item Build the grouped English database from PhrasIS, with one source-pair
record containing label, provenance, split metadata, and later all language
candidates.
\item Sample trial/dev/evaluation splits by source pair, never by ordered
instance, so that no reverse direction or language variant crosses splits.
\item Generate Spanish and Basque candidates from multiple translation systems,
starting with the tested \texttt{gpt-5.4} and \texttt{gpt-5.6-sol} workflow and
adding further providers if they improve non-collapsed coverage.
\item Run automatic diagnostics: target-side collapse, source-copy detection,
missing candidate detection, and later label-preservation checks after
acceptance.
\item Automatically reject collapsed candidates, not necessarily the whole
source pair.  If another model gives a non-collapsed candidate, that candidate
is retained for validation.
\item Send retained candidates to bilingual review.  Reviewers check phrase
adequacy, preservation of the intended NLI label, and preservation of the
deterministic reverse label.
\item Correct a candidate when a natural relation-preserving translation is
available; otherwise remove the language-specific item from official release.
\item Report final counts by source class, DiCo-NLI label, language, split,
generator, and rejection reason.
\end{enumerate}

This protocol directly addresses the reviewer concern.  The multilingual data
will not be accepted merely because an API produced fluent translations.  The
translation process is auditable, model-diverse, checkpointed, automatically
screened for structural failures, and followed by bilingual validation before
release.

\section{Current Trial Bundle}

The accompanying \texttt{trial\_data/} folder is generated from the accepted
non-collapsed candidates selected above.  It contains the grouped organizer JSON,
monolingual and mixed-language participant files, hidden gold files, submission
templates, translation-audit reports, and scored random/majority baselines.

\begin{table}[h]
\centering
\small
\begin{tabular}{lr}
\toprule
DiCo-NLI label & Trial pairs \\
\midrule
\texttt{BACKWARD\_ENTAILMENT} & 125 \\
\texttt{EQUIVALENCE} & 86 \\
\texttt{FORWARD\_ENTAILMENT} & 125 \\
\texttt{NEGATIVE\_OTHER} & 125 \\
\midrule
\textbf{Total accepted} & \textbf{461} \\
\textbf{Rejected by collapse filter} & \textbf{39} \\
\bottomrule
\end{tabular}
\caption{Current trial grouped JSON after non-collapsed candidate selection.}
\label{tab:trial-counts}
\end{table}

The accepted trial has zero Spanish or Basque exact-collapse pairs under the
automatic audit.  Candidate provenance is retained in the grouped JSON metadata:
Spanish uses 448 \texttt{gpt-5.4} candidates and 13 \texttt{gpt-5.6-sol}
fallback candidates; Basque uses 444 \texttt{gpt-5.4} candidates and 17
\texttt{gpt-5.6-sol} fallback candidates.

\begin{table}[h]
\centering
\small
\begin{tabular}{lrrrr}
\toprule
Track & Labeled & Unlabeled & Hidden gold & Template \\
\midrule
\texttt{track1} English & 797 & 797 & 797 & 797 \\
\texttt{track2} Spanish & 797 & 797 & 797 & 797 \\
\texttt{track3} Basque & 797 & 797 & 797 & 797 \\
\texttt{track4} Mixed & 4782 & 4782 & 4782 & 4782 \\
\bottomrule
\end{tabular}
\caption{Participant-facing files generated from the current trial bundle.}
\label{tab:trial-files}
\end{table}

\begin{table}[h]
\centering
\small
\begin{tabular}{llrrr}
\toprule
System & Track & Weighted F1 & SoftCons & HardCons \\
\midrule
Random & \texttt{track1} & 0.270 & 0.208 & 0.077 \\
Random & \texttt{track2} & 0.255 & 0.244 & 0.095 \\
Random & \texttt{track3} & 0.226 & 0.205 & 0.062 \\
Random & \texttt{track4} & 0.254 & 0.186 & 0.061 \\
\midrule
Majority & \texttt{track1} & 0.150 & 0.000 & 0.000 \\
Majority & \texttt{track2} & 0.150 & 0.000 & 0.000 \\
Majority & \texttt{track3} & 0.150 & 0.000 & 0.000 \\
Majority & \texttt{track4} & 0.150 & 0.000 & 0.000 \\
\bottomrule
\end{tabular}
\caption{Naive baselines recomputed on the current trial files.}
\label{tab:trial-baselines}
\end{table}

These files demonstrate the complete participant interface for the task while
using the same translation-validation policy planned for the full release.

\end{document}
